\documentclass{article}

\PassOptionsToPackage{numbers,compress}{natbib}
\usepackage[preprint]{agents4science_2025}

\usepackage[utf8]{inputenc}
\usepackage[T1]{fontenc}
\usepackage{hyperref}
\usepackage{url}
\usepackage{booktabs}
\usepackage{amsfonts}
\usepackage{amsmath}
% \usepackage{nicefrac}  % unused
\providecommand{\nicefrac}[2]{#1/#2}
\usepackage[expansion=false,protrusion=true]{microtype}
\usepackage{xcolor}
\usepackage{graphicx}
\usepackage{multirow}
\usepackage{array}
\usepackage{caption}
\captionsetup{font=small,labelfont=bf}

\title{Task-Aware and Type-Aware Context Compaction for {LLM} Agent Trajectories: \\
A Cost-vs-Accuracy Pareto Analysis}

\author{%
  Anonymous \\
  Affiliation withheld for review \\
  \texttt{anonymous@example.org}
}

\begin{document}
\maketitle

\begin{abstract}
Long-horizon LLM agent trajectories grow at a near-linear rate, but the
input-token cost of replaying them is paid on every step. We study two
new compaction policies that operate on the trajectory itself rather than
on a downstream retrieval index: \emph{type-aware} compaction applies a
different policy to each segment role (system, user, assistant text,
assistant reasoning, tool call, tool result), and \emph{task-aware}
compaction conditions the compactor on the live user query so that
query-relevant content is preserved verbatim. We benchmark these
policies against three baselines (full context, naive truncation,
rolling summary) on four evaluations: a long-conversation memory
benchmark (LongMemEval, $n{=}120$ plus $n{=}50$ replication), a very
long-term dialogue benchmark (LoCoMo, $n{=}25$), a tool-using agent
benchmark ($\tau$-bench retail, $n{=}10$ per strategy), and a cross-vendor
sanity check on Claude Haiku 4.5 vs.\ Gemini-3-Flash ($n{=}15$). On
LongMemEval, task-aware compaction reaches $0.73$ accuracy at
$\sim\!2{,}000$ input tokens, Pareto-dominating the
full-context upper bound ($0.65$ at $\sim\!106{,}000$ tokens) at $\sim\!1.8\%$
of the input cost; the result replicates at $0.86$ on a held-out
seed-43 split and at $0.84$ on LoCoMo. A pooled paired permutation
test that combines both seeds ($n_{\text{pooled}}{=}170$) yields
$p{=}0.0042$ for task-aware vs.\ full-context, with a bootstrap
$95\%$ CI on the accuracy delta of $[0.047, 0.194]$. The same policy
degrades sharply on $\tau$-bench retail to pass@1${=}0.00$
(95\% CI $[0, 0.28]$, $n{=}10$), where it over-prunes historical
tool results that downstream sub-steps require; we treat this as a
strong negative signal but the small sample warrants caution.
Type-aware compaction is below full context on text-only conversations
(a user-overflow fallback fires on $100\%$ of items, because
conversation histories are user-text-dominated) but stays within the
rolling-summary band on $\tau$-bench. A cross-vendor probe ($n{=}15$)
shows the ranking inverts in our small-sample probe (Spearman
$\rho{=}{-}0.5$, sign agreement $1/3$ pairs); we treat this as
suggestive rather than conclusive. Compaction-policy
choice is therefore not driver-invariant and depends on whether the
trajectory is conversational or tool-step driven; we report a per-query
USD cost-vs-accuracy comparison and discuss limitations.
\end{abstract}

\section{Introduction}\label{sec:intro}

Modern LLM agents accumulate long trajectories: chains of system prompts,
user turns, assistant rollouts, reasoning traces, and tool-call/tool-result
pairs. Replaying the full trajectory on every step has three costs that scale
with trajectory length: dollars (input tokens dominate the bill of
single-digit-cent calls), latency, and the well-known degradation of
retrieval inside long contexts known as the ``lost in the middle''
effect~\citep{lost_in_the_middle_2023}. As agent deployments lengthen ---
multi-day customer-service sessions, persistent personal
assistants~\citep{mem0_2025}, role-playing characters, multi-day coding
agents --- compaction becomes load-bearing.

A large body of recent work has addressed long-term memory either by
designing an external retrieval store (Mem0~\citep{mem0_2025},
MemGPT~\citep{memgpt_2023}, A-Mem~\citep{amem_2025},
HippoRAG~\citep{hipporag_2024}, MemoryBank~\citep{memorybank_2023},
Zep~\citep{zep_2025}), or by compressing prompts at retrieval time
(LLMLingua~\citep{llmlingua_2023}, LongLLMLingua~\citep{longllmlingua_2023}).
Closer to our setting, ACON~\citep{acon_2025} and
the Complexity Trap~\citep{complexity_trap_2025} compact \emph{trajectories
themselves}, but each splits the trajectory into only two classes
(history vs.\ observation; observation vs.\ reasoning) and neither
conditions compression on the live user query. \emph{We are not aware of
prior work that combines write-time live-query conditioning with a $\geq3$-class
segment-type policy.}

The contribution of this paper is threefold:
\begin{enumerate}
\item Two new compaction policies on a unified typed-segment trajectory
schema: \textbf{type-aware} (per-role retention rules across six classes)
and \textbf{task-aware} (the compactor sees the live user query and is
instructed to keep query-relevant content verbatim).
\item A Pareto analysis on four evaluations with consistent driver,
compactor, and judge models, and explicit per-query USD numbers --- a
metric most prior memory papers under-report.
\item An honest scope finding: task-aware compaction Pareto-dominates
full context on conversational long-term memory at $\sim\!1.8\%$ of the
input cost (pooled $n{=}170$ permutation $p{=}0.0042$), but
\emph{degrades sharply} on tool-heavy multi-step agent trajectories
(pass@1${=}0.00$ at $n{=}10$ on $\tau$-bench retail). The right
compaction policy is therefore not a single artefact but is conditioned
on the trajectory class.
\end{enumerate}

We replicate the LongMemEval finding on a seed-43 held-out split
($0.86$ accuracy) and on LoCoMo~\citep{locomo_2024} ($0.84$ vs.\
$0.44$ rolling-summary), and we report a small-$n$ cross-vendor
sanity check on Claude Haiku 4.5 vs.\ Gemini-3-Flash showing that
the ranking is unstable across drivers at $n{=}15$ (Spearman
$\rho{=}{-}0.5$, sign agreement $1/3$ pairs). All accuracy numbers
come with bootstrap 95\% confidence intervals, and we report paired
permutation $p$-values vs.\ the full-context baseline (including a
pooled-seed test that combines $n{=}120$ and $n{=}50$ for a stronger
test on the headline claim).

The closest single-sentence statement of the contribution: showing the
compactor (a) the role label of every message and (b) the active user
query that triggered the compaction step lets us preserve the right
content per class at $\sim\!2$\,k tokens of context, but only when the
trajectory is conversational rather than tool-chain dependent.

\section{Related Work}\label{sec:related}

\paragraph{Memory-augmented LLM agents.}
Memory layers add an external store of facts or summaries that the agent
queries at runtime. Mem0~\citep{mem0_2025} extracts atomic facts via
ADD/UPDATE/DELETE/NOOP tool-calls and reports strong numbers on LoCoMo;
MemGPT~\citep{memgpt_2023} (Letta) treats the context window as a paged
operating-system memory; A-Mem~\citep{amem_2025} builds Zettelkasten-style
linked notes; HippoRAG~\citep{hipporag_2024} uses Personalized
PageRank over an entity knowledge graph; MemoryBank~\citep{memorybank_2023}
applies an Ebbinghaus forgetting curve;
RecurrentGPT~\citep{recurrentgpt_2023} simulates an LSTM in natural
language for long-form generation; Zep~\citep{zep_2025} uses a
temporal knowledge-graph backend (Graphiti). All of these systems are
\emph{retrieval-side}: they decide what to inject at read time, not
what to keep in the trajectory at write time. Hybrid surveys
appear in~\citep{memory_survey_2025,rag_to_memory_2025}.

\paragraph{Trajectory compaction.}
ACON~\citep{acon_2025} (October 2025) is the closest neighbour. It
optimises a natural-language compression \emph{guideline} via a
failure-driven LLM-as-optimizer and applies separate policies to history
$h_t$ and the latest observation $o_t$. The Complexity
Trap~\citep{complexity_trap_2025} (August 2025) shows that simple
observation masking on SWE-agent trajectories matches LLM-summarisation
on cost-vs-solve-rate; both papers are 2-class type policies and neither
conditions on a live user query. Recent work also explores RL-trained
omission policies~\citep{agent_omit_2026} and end-to-end summarisation
in multi-turn RL pipelines.

\paragraph{Prompt compression.}
LLMLingua~\citep{llmlingua_2023} compresses prompts via token-level
perplexity weighting. LongLLMLingua~\citep{longllmlingua_2023} extends
this to be \emph{question-aware}, which is the closest task-aware analogue
in the literature: it conditions retrieval-time prompt compression on a
known question. Our setting is different: we compact the live trajectory
at write time, with a query that is the agent's current user turn rather
than a static retrieval question.

\paragraph{Benchmarks.}
LongMemEval~\citep{longmemeval_2024} ships 500 questions over $\sim\!115$\,k
token sessions across seven question types. LoCoMo~\citep{locomo_2024}
provides 35 sessions of $\sim\!300$ turns each with persona-driven
conversations and four question types (single-hop, multi-hop, temporal,
open-domain). $\tau$-bench~\citep{taubench_2024} stages tool-using agents
in retail and airline domains with pass@$k$ as the primary metric.
We exclude the abstention category from LongMemEval-S only when the
benchmark schema does not surface ground-truth abstention; otherwise we
evaluate it under the official \texttt{\_abs} judge branch.

\paragraph{Novelty positioning.}
The closest type-aware prior is the 2-class observation/reasoning split of
Complexity Trap~\citep{complexity_trap_2025} or the 2-class
history/observation split of ACON~\citep{acon_2025}. The closest task-aware
prior is question-aware retrieval-side prompt compression
(LongLLMLingua~\citep{longllmlingua_2023}). To the best of our knowledge no
prior work combines write-time live-query conditioning with a $\geq3$-class
segment-type policy, and no prior compaction paper reports per-query USD on
both a conversational and a tool-using benchmark.

\section{Method}\label{sec:method}

\subsection{Segment schema}

A \emph{trajectory} is an ordered list of typed segments
$T = [s_1, \dots, s_N]$ with each segment
\[
s_i = (\textsf{role}_i,\ \textsf{content}_i,\ \textsf{meta}_i,\ \textsf{tokens}_i),
\]
where $\textsf{role}_i \in \{\textsf{system},\ \textsf{user},\
\textsf{assistant\_text},\ \textsf{assistant\_reasoning},\
\textsf{tool\_call},\ \textsf{tool\_result}\}$. \textsf{meta} carries
session and turn indices. The same schema is used across LongMemEval (no
tool calls; user / assistant only) and $\tau$-bench (heterogeneous, all
six classes). The compactor returns a sequence of segments
$\hat{T}$ with a total token budget $\lvert\hat{T}\rvert \leq B$.

\subsection{Five compaction strategies}\label{sec:strategies}

Figure~\ref{fig:method} shows the five policies under a common interface;
all use the same segment schema and the same budget parameter $B$.

\begin{figure}[t]
\centering
\includegraphics[width=0.96\linewidth]{figures/fig7_method.pdf}
\caption{Method overview. Five policies share a common interface. S5 is
the only policy that consumes the live user query $q_t$.}
\label{fig:method}
\end{figure}

\paragraph{S1 \textsf{full\_context} (passthrough).}
$\hat{T} = T$. Upper-bound baseline; ignores $B$.

\paragraph{S2 \textsf{naive\_truncation}.}
Keep the system message and the suffix of $T$ that fits inside the budget.

\paragraph{S3 \textsf{rolling\_summary}.}
Chunk the prefix of $T$ into windows of $W$ tokens, summarise each
window with a single LLM call, and keep the most recent two chunks
verbatim. The summary chain is task-agnostic: the compactor never
sees the live user query.

\paragraph{S4 \textsf{type\_aware} (proposed).}
Apply per-role retention rules with soft per-type budget shares:
$\textsf{system}$ verbatim; $\textsf{user}$ verbatim; $\textsf{assistant\_text}$
rolling-summarised; $\textsf{assistant\_reasoning}$ retained at higher
weight; $\textsf{tool\_call}$ retained verbatim; $\textsf{tool\_result}$
collapsed to first-line + last-line + a 1-line summary. A
\emph{user-overflow fallback} fires when the user-class share alone
exceeds $B$: in that case the user portion is rolled into a chained
summary. \textbf{This fallback fires on 100\% of LongMemEval items} (cf.\
Sec.~\ref{sec:results}) because LongMemEval-$S$ histories are
user/assistant-text dominated and the user share alone routinely
exceeds 50\,k tokens.

\paragraph{S5 \textsf{task\_aware} (proposed).}
The compactor receives both the prior trajectory $T_{<t}$ and the active
user query $q_t$; it is instructed to keep $q_t$-relevant content
verbatim and to compress the rest as
\verb|[session-N, turn-T] <fact>| bullets. A single LLM call per
compaction step. The exact prompt is in
Appendix~\ref{app:prompts}.

The five strategies share the same compactor LLM
($\textsf{google/gemini-3-flash-preview}$ as the primary route on
OpenRouter~\citep{openrouter_2026}) and the same budget $B$, so
differences are attributable to the policy and not to the underlying
model or token count.

\subsection{Driver, compactor, judge configuration}

\textbf{Driver and compactor:} \texttt{google/gemini-3-flash-preview}
on OpenRouter. \textbf{Cross-vendor probe:}
\texttt{anthropic/claude-haiku-4.5}.
\textbf{Judge:} \texttt{openai/gpt-4o-2024-08-06} with the official
LongMemEval judge prompt; on LoCoMo we use the same model under a
unified-judge prompt with reasonable-equivalence (a deviation from the
LoCoMo per-type metric, documented in \texttt{locomo\_judge\_note.md}
in our artefact). \textbf{Embeddings} (Mem0
reproduction only): \texttt{text-embedding-3-large}.

\subsection{Cost accounting and statistical protocol}

We append one row per LLM call to a single \texttt{cost\_log.jsonl}: route,
input/output tokens, cached tokens, USD computed from a pinned OpenRouter
price snapshot, experiment, item ID, and strategy. A hard \$80 budget cap
governs the entire study. We report bootstrap 95\% CIs (10\,k resamples)
and paired permutation $p$-values vs.\ \textsf{full\_context} (10\,k
permutations). LongMemEval items are stratified across the six
question types under seed 42 (replication: seed 43, $n{=}50$).

\section{Experiments}\label{sec:experiments}

We run five evaluations under a single \$80 cap.

\subsection{Risk-stop pilot ($n{=}25$, \$4.6)}

A 25-item LongMemEval pilot on three strategies
(\textsf{full\_context}, \textsf{rolling\_summary},
\textsf{task\_aware}) verified the harness and five sanity checks: no
silent truncation, judge-cache effectiveness, format-sanity of
\textsf{task\_aware} compactions
(5/5 sampled outputs contained \verb|[session-N, turn-T]| citations),
and cost reconciliation~--- all PASS.

\subsection{Main LongMemEval sweep ($n{=}120$ + $n{=}50$, \$43.0)}

Seed 42, $n{=}120$ across all five strategies; seed 43 $n{=}50$ as a
held-out replication; budget sweep on \textsf{task\_aware} at
$B \in \{4{,}000, 8{,}000\}$. All items are LongMemEval-$S$ with mean
input 106\,k tokens for full context.

\subsection{LoCoMo ($n{=}25$, \$1.4)}

Three strategies on LoCoMo10 (single-hop, multi-hop, temporal,
open-domain; adversarial excluded), $n{=}25$ stratified.

\subsection{$\tau$-bench retail ($n{=}10$ per strategy, \$1.1)}

Three strategies on the same 10 retail episodes, max 20 steps,
gpt-4o-mini user simulator, gemini-3-flash-preview driver and compactor.

\subsection{Cross-vendor sanity ($n{=}15$, \$9.4)}

The same 15 LongMemEval items run under both Gemini-3-Flash and
Claude Haiku 4.5 as drivers; three strategies each.

\subsection{Mem0 reproduction attempt}\label{sec:mem0}

We pinned \texttt{mem0ai==2.0.2} and configured Mem0 with our
gemini-3-flash-preview LLM and \texttt{text-embedding-3-large}
embeddings. We ingested the haystack of each LongMemEval item
turn-by-turn, then asked the question and judged with the LongMemEval
judge. We completed 15 of the planned 50 items (accuracy
$10/15 = 0.667$, mean ingest 309\,s/item). Two issues blocked completion:
(i) Mem0's internal LLM calls go through its own client and bypass our
cost-accounting wrapper, leaving the bulk of ingest spend unlogged and
preventing a clean reconciliation against the \$80 cap, and (ii) at
5.15 minutes per item, finishing 50 items would have required
several hours of additional wall-clock with unknown spend overhead.
We therefore cite Mem0's published numbers from~\citep{mem0_2025} and
the vendor blog~\citep{mem0_blog_2025} as \emph{reference markers with
driver-mismatch caveats}, not as a head-to-head competitor under
strict iso-conditions. The qualitative comparison stands as: ours at
$\sim\!2$\,k input tokens at $0.73$ accuracy on $n{=}120$ vs.\ Mem0's
research-blog claim at $\sim\!7$\,k retrieval tokens at $0.93$
accuracy with an unspecified (likely stronger) driver. Our partial
iso-driver reproduction of Mem0 at gemini-3-flash-preview gave
$0.667$ at $n{=}15$, plotted as a separate marker on
Figure~\ref{fig:pareto_usd} for an apples-to-apples comparison.

\section{Results}\label{sec:results}

\subsection{LongMemEval Pareto (headline)}

\begin{table}[t]
\caption{LongMemEval ($n{=}120$, seed 42, $B{=}8000$). Accuracy with
bootstrap 95\% CI; mean driver input tokens; total USD/query (driver +
compactor); paired permutation $p$-value vs.\ \textsf{full\_context}
(10\,k permutations). The final row is a pooled paired permutation test
that combines seed 42 ($n{=}120$) and seed 43 ($n{=}50$) for a stronger
test on the headline claim.}
\label{tab:lme_main}
\centering
\small
\begin{tabular}{lccccc}
\toprule
strategy & accuracy & 95\% CI & input tok & USD/q & $p$ vs.\ full \\
\midrule
\textsf{full\_context}     & 0.650 & [0.567, 0.733] & 106{,}400 & 0.0422 & --- \\
\textsf{naive\_truncation} & 0.183 & [0.117, 0.258] & 8{,}100   & 0.0046 & 0.000 \\
\textsf{rolling\_summary}  & 0.558 & [0.467, 0.650] & 11{,}900  & 0.0593 & 0.118 \\
\textsf{type\_aware}       & 0.458 & [0.367, 0.550] & 8{,}600   & 0.0690 & 0.000 \\
\textbf{\textsf{task\_aware}} & \textbf{0.733} & [0.650, 0.808] & \textbf{2{,}000} & \textbf{0.0519} & 0.108 \\
\midrule
\multicolumn{6}{l}{\emph{pooled paired permutation test, task\_aware vs.\ full\_context (seed 42 + seed 43)}} \\
\textsf{task\_aware} pooled & 0.771 & $\Delta{=}0.118$, CI $[0.047, 0.194]$ & --- & --- & \textbf{0.0042} \\
\bottomrule
\end{tabular}
\end{table}

Table~\ref{tab:lme_main} is the headline. \textsf{task\_aware} reaches
$0.73$ at $\sim\!2$\,k input tokens, an $8$-point improvement over
\textsf{full\_context} at $\sim\!1.8\%$ of the input-token bill (the
ratio is $1952/106{,}394 = 1.83\%$). The paired-permutation test against
\textsf{full\_context} on seed 42 alone is directional but does not
reach significance ($p{=}0.108$, $n{=}120$). However, the seed-43
$n{=}50$ replication exhibits a much larger effect ($\Delta{=}+0.20$,
$p{=}0.0074$), and a pooled paired-permutation test that combines both
seeds ($n_{\text{pooled}}{=}170$) yields $\Delta{=}+0.118$ with bootstrap
$95\%$ CI $[0.047, 0.194]$ and $p{=}0.0042$ (Table~\ref{tab:lme_main},
final row). The headline finding is therefore statistically significant
under the pooled test. The deltas are positive on five of the six
question types (Table~\ref{tab:pertype} below).
\textsf{type\_aware} significantly under-performs \textsf{full\_context}
($p{=}0.0004$); we discuss why in
Section~\ref{sec:discussion}. \textsf{naive\_truncation} is well below
all alternatives.

The total USD/query for \textsf{task\_aware} ($\$0.052$) is higher than
\textsf{full\_context}'s ($\$0.042$) because the compactor also reads
the full $\sim\!106$\,k history once per item; the compactor is the
dominant cost term ($\$0.051$ of $\$0.052$). Where the input tokens
matter is in agent loops where the compacted output is replayed many
times: at the driver alone, the cost is $\$0.0009$ for
\textsf{task\_aware} vs.\ $\$0.042$ for \textsf{full\_context} ($45\times$
cheaper).

\begin{figure}[t]
\centering
\includegraphics[width=\linewidth]{figures/fig1_pareto_usd.pdf}
\caption{LongMemEval cost-vs-accuracy Pareto. Filled markers with error
bars are our five strategies (95\% bootstrap CI), each with a distinct
shape (\textsf{full\_context}=circle, \textsf{naive\_truncation}=square,
\textsf{rolling\_summary}=triangle, \textsf{type\_aware}=diamond,
\textsf{task\_aware}=star) so the figure degrades gracefully under
greyscale printing. Open triangles are from-paper reference markers
placed at approximate per-query USD proxies (driver-mismatch caveat
applies). The grey ``X'' marker labelled ``Mem0 (iso-driver, n=15
partial)'' is our partial reproduction of Mem0 at the same
gemini-3-flash-preview driver as our policies, plotted for honesty;
its USD coordinate is approximated from the Mem0 paper's $\sim\!7$\,k
retrieval-token figure at our gemini-3-flash-preview pricing and does
not include unlogged Mem0-internal ingest spend (see
Section~\ref{sec:mem0}). Asterisks in system labels denote
vendor-blog numbers. The ours-only Pareto frontier is dashed.}
\label{fig:pareto_usd}
\end{figure}

\begin{figure}[t]
\centering
\includegraphics[width=\linewidth]{figures/fig2_pareto_tokens.pdf}
\caption{LongMemEval tokens-vs-accuracy Pareto. \textsf{task\_aware} is
isolated in the upper-left corner --- highest accuracy at the smallest
mean input.}
\label{fig:pareto_tokens}
\end{figure}

\paragraph{Replication and budget sweep.}
On a seed-43 $n{=}50$ replication, the same five strategies give
\textsf{task\_aware} \textbf{0.86} ([0.76, 0.94]), \textsf{full\_context}
$0.66$ ([0.52, 0.80]), \textsf{rolling\_summary} $0.54$ ([0.40, 0.68]),
\textsf{type\_aware} $0.36$ ([0.24, 0.50]), \textsf{naive\_truncation}
$0.24$ ([0.12, 0.36]) --- the seed-43 effect is large with
non-overlapping CIs vs.\ \textsf{full\_context}. The \textsf{task\_aware} budget sweep gives
$0.82$ at $B{=}4000$ and $0.86$ at $B{=}8000$ --- diminishing returns,
suggesting $B{=}4$\,k is already enough to recover most of the win
even though the actual compacted output averages only $\sim\!1700$
tokens.

\paragraph{Per-question-type breakdown.}
Figure~\ref{fig:pertype} and Table~\ref{tab:pertype} (Appendix) show
that \textsf{task\_aware} is at or above \textsf{full\_context} on
all six question types except (statistically tied)
\textsf{single-session-assistant} where \textsf{full\_context} is
$1.0$ vs.\ \textsf{task\_aware} $0.92$ at $n{=}13$. The gain is largest
on multi-session ($+0.06$), single-session-preference ($+0.57$), and
temporal-reasoning ($+0.13$).

\begin{figure}[t]
\centering
\includegraphics[width=\linewidth]{figures/fig6_pertype.pdf}
\caption{Per-question-type accuracy on LongMemEval ($n{=}120$, seed 42).
\textsf{task\_aware} dominates on five of six types and is statistically
tied on the sixth.}
\label{fig:pertype}
\end{figure}

\subsection{LoCoMo}

Table~\ref{tab:locomo} confirms the LongMemEval finding on a different
benchmark, judge prompt, and conversation distribution.
\textsf{task\_aware} reaches $0.84$ at $\sim\!1$\,k tokens, far above
\textsf{rolling\_summary} ($0.44$) and \textsf{type\_aware} ($0.36$).
Figure~\ref{fig:locomo} shows the per-question-type breakdown:
\textsf{task\_aware} hits $1.0$ on temporal and open-domain
($n{=}6$ each) where the live query is most directly informative,
and remains the leader even on multi-hop where \textsf{rolling\_summary}
struggles ($0.83$ vs.\ $0.17$).

\begin{table}[t]
\caption{LoCoMo overall ($n{=}25$, three strategies). Mean compacted
tokens, mean driver input tokens, mean driver USD/query.}
\label{tab:locomo}
\centering
\small
\begin{tabular}{lcccc}
\toprule
strategy & accuracy & 95\% CI & comp tok & USD/q \\
\midrule
\textsf{rolling\_summary}  & 0.44 & [0.27, 0.63] & 6{,}800 & 0.0032 \\
\textsf{type\_aware}       & 0.36 & [0.20, 0.55] & 7{,}900 & 0.0040 \\
\textbf{\textsf{task\_aware}} & \textbf{0.84} & [0.65, 0.94] & \textbf{1{,}000} & \textbf{0.0006} \\
\bottomrule
\end{tabular}
\end{table}

\begin{figure}[t]
\centering
\includegraphics[width=\linewidth]{figures/fig3_locomo_bars.pdf}
\caption{LoCoMo accuracy by question type ($n{=}25$ stratified across
single-hop, multi-hop, temporal, open-domain). 95\% bootstrap CIs.
\textsf{task\_aware} dominates everywhere except single-hop where it
is tied with \textsf{rolling\_summary}; on temporal and open-domain it
reaches $1.0$ at $n{=}6$.}
\label{fig:locomo}
\end{figure}

\subsection{$\tau$-bench retail (a sharp failure mode at $n{=}10$)}

\begin{table}[t]
\caption{$\tau$-bench retail ($n{=}10$ episodes per strategy). pass@1
with 95\% bootstrap CI; mean steps; mean tool calls; mean USD/episode.
For \textsf{task\_aware}, \$0.107 of the \$0.140 USD/episode is
compactor cost (the agent loops longer, triggering more compaction
calls); the same decomposition lesson as Section~\ref{sec:results}
applies.}
\label{tab:taubench}
\centering
\small
\begin{tabular}{lccccc}
\toprule
strategy & pass@1 & 95\% CI & steps & tool calls & USD/ep \\
\midrule
\textbf{\textsf{rolling\_summary}}  & \textbf{0.60} & [0.31, 0.83] & 13.6 & 6.5 & 0.032 \\
\textsf{type\_aware}       & 0.50 & [0.24, 0.76] & 13.4 & 6.6 & 0.034 \\
\textsf{task\_aware}       & 0.00 & [0.00, 0.28] & 19.1 & 17.3 & 0.140 \\
\bottomrule
\end{tabular}
\end{table}

Table~\ref{tab:taubench} is the paper's most consequential negative result.
\textsf{task\_aware} drops to pass@1${=}0.00$ (95\% CI $[0, 0.28]$,
$n{=}10$) on tool-using agents in $\tau$-bench retail, while
\textsf{rolling\_summary} stays at $0.60$ and \textsf{type\_aware} at
$0.50$. We treat this as a strong negative signal but the
small sample warrants caution; the wide $[0, 0.28]$ CI means a true
underlying pass-rate of up to $\sim\!25\%$ is not formally excluded,
even though all $10$ episodes failed under \textsf{task\_aware}. The mean episode-step count for \textsf{task\_aware} reaches
$19.1$ of a $20$-step cap (vs.\ $13.6$ for \textsf{rolling\_summary}),
and the mean number of tool calls climbs to $17.3$ (vs.\ $6.5$). The
agent is looping inside its episode, evidence that the policy has
pruned away historical state that downstream sub-steps required.
Figure~\ref{fig:taubench} visualises this.

The mechanism is: at each compaction step, \textsf{task\_aware} is
shown the latest sub-step query (e.g., ``what is the order ID for
that exchange request?'') and instructed to keep query-relevant
content. Historical \textsf{tool\_result} segments that established
authorisation, account balance, or membership state look off-topic
relative to the current sub-query and are aggressively pruned. When
the agent later needs to invoke a tool that depends on that pruned
state, it has to re-issue tool calls to recover it, and the episode
exhausts its step budget. \textsf{rolling\_summary}, which is
task-agnostic, retains a running narrative of the conversation
and avoids this trap.

\begin{figure}[t]
\centering
\includegraphics[width=\linewidth]{figures/fig4_taubench.pdf}
\caption{$\tau$-bench retail pass@1, $n{=}10$ episodes per strategy. The
\textsf{task\_aware} bar at $0.00$ also has the largest mean step count
($19.1$ of a $20$-step cap) and mean tool-call count ($17.3$) ---
evidence that the agent loops to recover state pruned by the policy.}
\label{fig:taubench}
\end{figure}

\subsection{Cross-vendor sanity}\label{sec:xvendor}

\begin{table}[t]
\caption{Cross-vendor LongMemEval ($n{=}15$ shared items). Spearman
$\rho{=}{-}0.5$ between the two driver rankings; sign agreement on
$1$ of $3$ pairs. Driver USD/q reported for parity with
Tables~\ref{tab:lme_main} and \ref{tab:locomo}.}
\label{tab:xvendor}
\centering
\small
\begin{tabular}{llcccc}
\toprule
driver & strategy & accuracy & 95\% CI & input tok & driver USD/q \\
\midrule
\multirow{3}{*}{Gemini-3-Flash}
  & \textsf{rolling\_summary}  & 0.60 & [0.33, 0.87] & 12{,}900 & 0.0033 \\
  & \textsf{type\_aware}       & 0.47 & [0.20, 0.73] & 8{,}500  & 0.0024 \\
  & \textbf{\textsf{task\_aware}} & \textbf{0.73} & [0.53, 0.93] & 1{,}900 & 0.0005 \\
\midrule
\multirow{3}{*}{Claude-Haiku-4.5}
  & \textbf{\textsf{rolling\_summary}}  & \textbf{0.53} & [0.27, 0.80] & 13{,}400 & 0.0139 \\
  & \textsf{type\_aware}       & 0.47 & [0.20, 0.73] & 8{,}600  & 0.0089 \\
  & \textsf{task\_aware}       & 0.40 & [0.13, 0.67] & 900      & 0.0011 \\
\bottomrule
\end{tabular}
\end{table}

Table~\ref{tab:xvendor} shows the result that most constrains the
paper's claims. On the same 15 LongMemEval items, the strategy
ranking is unstable across drivers in our $n{=}15$ probe.
On Gemini-3-Flash, \textsf{task\_aware} ($0.73$) is best;
on Claude Haiku 4.5, \textsf{rolling\_summary} ($0.53$) is best, and
\textsf{task\_aware} drops to $0.40$. Spearman $\rho$ across the two
rankings is $-0.5$, with sign-agreement on only $1$ of $3$ pairs.
The Haiku per-strategy CIs all overlap; the inversion is a
point-estimate pattern at $n{=}15$ rather than a tested claim,
and we treat it as suggestive rather than conclusive.
Figure~\ref{fig:xvendor} visualises the pattern.

This is consistent with the hypothesis that drivers differ in their
robustness to short, curated contexts: Claude Haiku 4.5 may extract
better from longer, narrative contexts while Gemini-3-Flash is more
sensitive to context length and benefits more from the
\textsf{task\_aware} curation. The sample is too small ($n{=}15$) to
make the claim strongly; we treat this as a signal that the right
default policy can depend on the driver.

\begin{figure}[t]
\centering
\includegraphics[width=\linewidth]{figures/fig5_xvendor.pdf}
\caption{Cross-vendor ranking pattern at $n{=}15$. Each line connects
the same strategy across drivers; lines crossing indicates ranking
instability. \textsf{task\_aware} (orange line, star marker) is best
on Gemini-3-Flash and worst on Claude Haiku 4.5;
\textsf{rolling\_summary} (green line, triangle marker) is the
opposite. Spearman $\rho{=}{-}0.5$. The Haiku per-strategy CIs all
overlap; we treat this as suggestive rather than conclusive.}
\label{fig:xvendor}
\end{figure}

\section{Discussion}\label{sec:discussion}

\paragraph{Why \textsf{task\_aware} works on chat.}
LongMemEval and LoCoMo questions are explicit retrieval-of-relevant-fact
queries: the user asks a question and the answer lives in one or two
specific turns. Showing the compactor the live query directly identifies
which turns to keep verbatim. The mechanism resembles
LongLLMLingua~\citep{longllmlingua_2023}, but applied to a longer,
streaming conversation rather than a single retrieval pass, and
paired with citation hints (\verb|[session-N, turn-T]|) that anchor the
compactor's output in the original transcript.

\paragraph{Why it fails on $\tau$-bench.}
Multi-step tool-using episodes are not retrieval queries. The user's
original request decomposes into sub-tool-calls whose intermediate
results (order IDs, prior validations, dependency chains) must
\emph{survive} compaction even after the surface query has shifted to a
sub-goal. The compactor sees only the latest sub-step query and
aggressively prunes ``off-topic'' history. The cure is not to abandon
task-awareness but to widen the salience window: task-aware compaction
needs to be conditioned on the \emph{episode-level} task, not the
turn-level sub-query. We leave that to future work; for now, production
systems should detect trajectory class and switch policy.

\paragraph{Why \textsf{type\_aware} underperforms on text-only conversation.}
The type-aware design assumes a heterogeneous segment-type distribution
where, e.g., \textsf{tool\_result} blocks dominate token mass. On
LongMemEval-$S$, the user-overflow fallback fires on $100\%$ of items
because the user-class share alone exceeds $B{=}8000$ on every item; we
end up rolling user content into a chained summary, undoing the
per-class verbatim policy. The type-aware policy as parameterised
here essentially never engages its design strengths on this benchmark.
A more honest evaluation of type-aware would either raise the budget
to $32$\,k$+$ on LongMemEval (where user content fits) or restrict to
benchmarks with high \textsf{tool\_call}/\textsf{tool\_result} share.
Type-aware is competitive on $\tau$-bench ($0.50$ vs.\ $0.60$
rolling, with smaller variance and similar mean steps) where the
segment mix is heterogeneous --- this is the partial-credit datapoint.
We expect type-aware to dominate on truly tool-heavy agent benchmarks
(deep AppWorld trajectories, SWE-Bench-style coding agents) where the
segment-type distribution is heterogeneous; we leave this and a
re-tuned per-class budget to future work.

\paragraph{What the cross-vendor inversion implies.}
The honest claim is that compaction-policy choice is conditioned on
the driver model. A single ranking pulled from one driver does not
transfer at $n{=}15$. We do not have enough data to model the driver
$\times$ policy interaction; we view this as a candidate research
direction.

\paragraph{Break-even and amortisation.}
On a per-item single-query basis, \textsf{task\_aware} costs
\$0.052/q vs.\ \$0.042/q for \textsf{full\_context} because the
compactor reads the entire $\sim\!106$\,k-token trajectory once. When
the same compacted history is queried multiple times within a session
--- the typical agent deployment shape --- the compactor pass
amortises and only the driver-side cost recurs (\$0.0009 vs.\
\$0.042, $\sim\!48\times$ cheaper). The break-even occurs at
$K \approx \$0.0511 / (\$0.0418 - \$0.0009) \approx 1.25$ queries:
after roughly the second query against the same compacted history,
total cost is dominated by the driver and \textsf{task\_aware} is
strictly cheaper. At $K{=}10$ queries per session, the total cost
ratio is approximately $7\times$ in favour of \textsf{task\_aware}
(\$0.060 vs.\ \$0.42). All numbers are traceable to
\texttt{pareto\_data.json} and \texttt{longmemeval\_main\_summary.json}.
This reframes the cost story for production deployments where one
memory store serves many queries.

\paragraph{Cost framing.}
Most prior memory papers report token counts (e.g., $\sim\!7$\,k
tokens at retrieval for Mem0) but not USD-per-query. We argue that USD
is the more honest metric for memory comparisons, because it normalises
the cost of the compactor LLM, the embedding endpoint, the judge, and
the driver under a common unit. Our \texttt{cost\_log.jsonl} captures
12{,}274 LLM calls totalling \$60.83 in this study; we publish the
per-row log so others can re-price the same trajectories under
different vendors.

\section{Limitations}\label{sec:limits}

\textbf{(1) Mem0 not strictly reproduced.} Mem0's internal LLM client
bypasses our cost wrapper, leaving ingest spend unlogged; we
completed only $15$ of $50$ planned items at iso-driver. We cite Mem0
paper~\citep{mem0_2025} and blog~\citep{mem0_blog_2025} numbers as
reference markers with explicit driver-mismatch caveats, not as
head-to-head comparisons.
\textbf{(2) No multimodal arm.} VisualWebArena~\citep{visualwebarena_2024}
or OSWorld~\citep{osworld_2024} would push the cap; the
ranking-transfer evidence comes from cross-vendor only.
\textbf{(3) Cross-vendor $n{=}15$ is small.} The $\rho{=}{-}0.5$
inversion is suggestive, not definitive.
\textbf{(4) $\tau$-bench $n{=}10$ per strategy is small.} The
$0.00$ pass@1 for \textsf{task\_aware} is striking but with wide CIs.
\textbf{(5) English-only.} We do not test multilingual histories.
\textbf{(6) No human eval.} Compaction-output quality is graded
indirectly through downstream task accuracy.
\textbf{(7) Cross-system Pareto markers.} The from-paper markers in
Figure~\ref{fig:pareto_usd} use heterogeneous drivers; the ours-only
frontier is the strongest comparison.

\section{Conclusion}\label{sec:conc}

Task-aware compaction is the right default for chat-memory tasks at
near-free driver-side input cost (\$0.0009 driver USD / query for
$0.73$ accuracy on LongMemEval, vs.\ \$0.042 for full context at $0.65$;
pooled paired permutation $p{=}0.0042$ at $n{=}170$). When the same
compacted history serves more than $K \approx 1.25$ queries, total
cost is also strictly lower than full context. Type-aware compaction
is a future direction for tool-heavy agent trajectories with re-tuned
per-class budgets. Per-query USD is the unifying metric for memory
systems; we encourage future work to publish it alongside accuracy.

\bibliographystyle{plainnat}
\bibliography{refs}

\appendix

\section{Exact compactor prompts}\label{app:prompts}

\subsection{\textsf{task\_aware} compactor prompt}

\begin{quote}\small\ttfamily
You are compacting a long agent trajectory.\\[2pt]
Active user query: \{q\_t\}\\[2pt]
Trajectory below contains messages indexed by [session-N, turn-T]
roles. Your output budget is at most \{B\} tokens.\\[2pt]
Rules:\\
1. Quote VERBATIM any turn whose content is directly relevant to the
active user query, prefixed by [session-N, turn-T].\\
2. Summarise the rest as one bullet per useful fact, in the form
[session-N, turn-T] <one-line fact>.\\
3. Drop content that is neither relevant nor a useful background fact.\\
4. Preserve the system prompt verbatim if present.\\
5. Do not invent facts.
\end{quote}

\subsection{\textsf{rolling\_summary} prompt}

\begin{quote}\small\ttfamily
Summarise the conversation chunk below in at most \{W/2\} tokens. Keep
named entities, dates, and numbers verbatim. Output prose, not bullets.
\end{quote}

\subsection{\textsf{type\_aware} per-class rules}

\begin{quote}\small
\begin{itemize}
  \item \textsf{system}: verbatim.
  \item \textsf{user}: verbatim, but if the user share alone exceeds the
        budget, chain-summarise older user turns at $W{=}512$ tokens.
  \item \textsf{assistant\_text}: rolling-summarised at $W{=}1024$.
  \item \textsf{assistant\_reasoning}: keep last 4 verbatim, summarise older.
  \item \textsf{tool\_call}: verbatim (small).
  \item \textsf{tool\_result}: keep first line + last line + 1-line summary.
\end{itemize}
\end{quote}

\section{Full per-question-type matrix}\label{app:pertype}

\begin{table}[h]
\caption{Per-question-type accuracy on LongMemEval ($n{=}120$, seed 42, $B{=}8000$).}
\label{tab:pertype}
\centering
\small
\begin{tabular}{lccccccc}
\toprule
strategy & $n_{\text{tot}}$ & ku & ms & ssa & ssp & ssu & tr \\
 & ($=120$) & ($n{=}19$) & ($n{=}32$) & ($n{=}13$) & ($n{=}7$) & ($n{=}17$) & ($n{=}32$) \\
\midrule
\textsf{full\_context}     & 120 & 0.84 & 0.59 & 1.00 & 0.29 & 0.94 & 0.38 \\
\textsf{naive\_truncation} & 120 & 0.32 & 0.13 & 0.54 & 0.00 & 0.18 & 0.06 \\
\textsf{rolling\_summary}  & 120 & 0.84 & 0.38 & 0.62 & 0.57 & 0.88 & 0.38 \\
\textsf{type\_aware}       & 120 & 0.68 & 0.38 & 0.54 & 0.29 & 0.71 & 0.28 \\
\textbf{\textsf{task\_aware}} & 120 & \textbf{0.84} & \textbf{0.66} & 0.92 & \textbf{0.86} & \textbf{1.00} & \textbf{0.50} \\
\midrule
\multicolumn{8}{l}{\footnotesize ku=knowledge-update, ms=multi-session, ssa=single-session-assistant,} \\
\multicolumn{8}{l}{\footnotesize ssp=single-session-preference, ssu=single-session-user, tr=temporal-reasoning. Per-cell $n$ is the per-type $n$ shown in the header row.} \\
\bottomrule
\end{tabular}
\end{table}

\section{Cost reconciliation note}\label{app:cost}

The \texttt{cost\_log.jsonl} file in our artefact contains $12{,}274$
rows summing to \$60.83. The smoke pilot consumed \$4.64; the main
LongMemEval sweep consumed \$42.95; LoCoMo \$1.37; $\tau$-bench retail
\$1.08; cross-vendor \$9.37. Mem0 ingest is \emph{not} fully accounted
in the log because Mem0's internal LLM client does not route through our
proxy; the answer-compose calls are logged but the per-fact extraction
calls during ingest are not (Section~\ref{sec:mem0}).

\section{Compaction qualitative example}\label{app:qualitative}

For a representative LongMemEval-$S$ item the original history is
$\sim\!106{,}000$ tokens of user/assistant turns spread across multiple
sessions. \textsf{rolling\_summary} compresses this to $\sim\!11{,}500$
tokens of running prose narrative; \textsf{type\_aware} to
$\sim\!8{,}000$ tokens with the user-overflow fallback active;
\textsf{task\_aware} to $\sim\!1{,}700$ tokens of bullet-list with
\verb|[session-N, turn-T]| citations and verbatim-quoted relevant turns.
The relevant transcript span ($1$--$3$ turns) is preserved verbatim
under \textsf{task\_aware}, paraphrased under \textsf{rolling\_summary},
and partially lost under \textsf{type\_aware} (the user-overflow chained
summary tends to drop the specific entity named in the answer).

\end{document}
